% !PRINTED 203
\weberpage{203}
%<<CONT
and we obtain a second possibility:

\gesperrt{II. Dihedral group,} $\nu = \nu' = 2, \nu'' = m, n = 2m, m \geqq 2$

\[ \mu = \mu' = m, \quad \mu'' = 2. \]

If, further, neither of the numbers $\nu', \nu''$ is equal to 2, then one of them must be equal to 3. For if they are both $\geqq 4$, then

\[ \mu' + \mu'' \leqq \frac{n}{2}, \]

which stands in contradiction with (16). If then $\nu' = 3$, $\mu' = \dfrac{n}{3}$, it follows from (16):

\weq{17}{\mu'' = \frac{n}{6} + 2, \quad \nu'' = \frac{6n}{n+12},}

from which it follows that we must have $\nu'' < 6$, and hence can have only one of the values 3, 4, 5.

Accordingly we obtain three more possible cases:

\gesperrt{III. Tetrahedral group:} $\nu = 2, \nu' = 3, \nu'' = 3, n = 12$

\[ \mu = 6, \quad \mu' = 4, \quad \mu'' = 4. \]

\gesperrt{IV. Octahedral group:} $\nu = 2, \nu' = 3, \nu'' = 4, n = 24$

\[ \mu = 12, \quad \mu' = 8, \quad \mu'' = 6. \]

\gesperrt{V. Icosahedral group:} $\nu = 2, \nu' = 3, \nu'' = 5, n = 60$

\[ \mu = 30, \quad \mu' = 20, \quad \mu'' = 12. \]

With this all possibilities are exhausted.

It is, to be sure, not yet proved that these groups, which we have provisionally listed under their customary names, and which we gather under the common name of the \gesperrt{polyhedral groups}, actually exist, nor how great their manifoldness is. Only the following considerations lead to this proof.

\begin{center}
\S.\,53.

\gesperrt{Transformation of the substitutions of $G$ to simple forms.}
\end{center}

Before we proceed to the definitive establishment of these groups, we derive a theorem which restricts still further the possibility of the substitutions that occur.
